\documentclass[10pt,a4paper]{amsart}
\usepackage[margin=19mm]{geometry}
\usepackage{amsmath,amssymb,mathtools}
\usepackage[scr=rsfs,cal=euler]{mathalfa}
\usepackage{xfrac}
\usepackage[unicode,hidelinks,bookmarks=false]{hyperref}
\newcommand{\ie}{i.e.\ }
\newcommand{\cf}{cf.\ }
\newcommand{\resp}{respectively\ }
\newcommand{\Subsection}{\subsection}
\providecommand{\nobreakdash}{\nobreak\hbox{-}}
\setlength{\emergencystretch}{2em}
\multlinegap0pt
\usepackage{amssymb}
\usepackage{xcolor}
\newtheorem{theorem}{Theorem}
\title{The Lingering Phenomenon and Pattern Formation in a Nonlocal Population Model with Cognitive Map}
\author{Kyung-Han Choi, Thomas Hillen}
\date{Source: arXiv:2602.16022v1 (CC BY 4.0)}
\begin{document}
\maketitle

\noindent\emph{Source and licence.} The Lingering Phenomenon and Pattern Formation in a Nonlocal Population Model with Cognitive Map. Version arXiv:2602.16022v1 (CC BY 4.0), distributed by arXiv under the Creative Commons Attribution 4.0 International licence, http://creativecommons.org/licenses/by/4.0/, as recorded in the arXiv licence metadata for this exact version and shown on the abstract page https://arxiv.org/abs/2602.16022v1. Full licence notice: \texttt{licenses/arxiv-2602.16022-CC-BY-4.0.txt}.\par\smallskip

\noindent\textit{Editorial scope.} Counted unit: the article's theorem thm:lingering (source0.tex lines 289--307). The article's complete original proof of it is delivered with it (source0.tex lines 309--434, one \texttt{proof} environment of 126 lines). No mathematics is written, repaired or re-derived: the statement and the proof are the article's own lines, in the article's own notation, and no retained line is substituted or re-flowed.  Retained uncounted as the source-local context the proof needs: lines 108--116; lines 149--160; lines 235--246; lines 248--251.  Nothing was omitted from any retained span, so the package carries no omission record.  No retained line contains a citation, so the wrapper carries no bibliography and no bibliography entry.  No retained line contains a figure environment, an image-inclusion command or drawing code, so no image is delivered and the package declares no asset dependency.  Editorial formatting only, in addition: an AMS \texttt{amsart} preamble that restates the article's own declarations and macros used by the retained lines, namely the environments \texttt{theorem} on separate counters, together with the standard packages those lines need.  The retained environments print in this document as Theorem 1 in order of appearance, whereas the article prints the same items with the numbering of its own full document.  The complete native source of the article as distributed in the arXiv e-print for this exact version is bundled unchanged as \texttt{sources/194842/source0.tex}.
\begin{equation}\label{defkernel}K(x,y) = \left\{\begin{aligned}
    &J_R(x-y) \text{ or }\frac{J_R(x-y)}{Z_R(x)} &&\text{ if } x-y\in B_R(0),\\
    &\qquad\qquad 0&& \text{ otherwise, }
\end{aligned}\right. 
% \text{ and }  \chi_{\overline{\Omega}}(x) =\left\{\begin{aligned}
%     &1 && \text{ if }x\in\overline\Omega\\
%     &0&& \text{ otherwise. }
% \end{aligned}\right.
\end{equation}

\begin{equation}\label{eqn: dispersalonly}
\left\{
\begin{aligned}
&\partial_t u  = \Delta\big(\gamma(m) u\big)
&& \text{in } \Omega\times(0,T], \\[4pt]
&\partial_t m = \big(\alpha\,\bar s(x)-m\big)\,u - \mu\,m
&& \text{in } \Omega\times(0,T]\\
&\nabla \big(\gamma(m) u\big)\cdot \mathbf{n} = 0
&& \text{on } \partial\Omega\times(0,T],
\end{aligned}
\right.
\end{equation} where \(\gamma:[0,\infty)\to(0,\infty)\) is a positive, strictly decreasing motility function. The perceived quantity \(\bar s(x)\) is defined on \(\Omega\) as in~\eqref{defkernel} and is independent of time in this section. Throughout the simulations below, we choose

\begin{equation}\label{eqn: SS dispersal}
\left\{
\begin{aligned}
&0  = \big(\gamma(m_\infty) u_\infty\big)''
&& \text{in } (-\ell,\ell)\times(0,\infty), \\[4pt]
&0 = \big(\alpha\,\bar s(x)-m_\infty\big)\,u_\infty - \mu\,m_\infty
&& \text{in } (-\ell,\ell)\times(0,\infty),\\
&\big(\gamma(m_\infty) u_\infty\big)'(\pm\ell) = 0
&& \text{on } (0,\infty)
\end{aligned}
\right.
\end{equation}

$u_\infty=u_\infty(x;\alpha,\mu)$ and $m_\infty=m_\infty(x;\alpha,\mu)$. Since the homogeneous Neumann boundary condition implies conservation of total mass for \eqref{eqn: dispersalonly}, the limiting steady state satisfies that \begin{equation}
    \label{mass conservation}
    \int_{-\ell}^\ell u_\infty(x;\alpha,\mu) dx = M,
\end{equation} where $M>0$ is determined by the initial mass of \eqref{eqn: dispersalonly} and hence is independent of $\alpha$ and $\mu$.

\begin{theorem}[Lingering phenomenon on unimodal $\bar s(x)$]\label{thm:lingering}
Let $\gamma\in C^2[0,\infty)$. Assume that $\bar s(x)$ is even in $(-\ell,\ell)$ and decreasing in $(0,\ell)$. Given $\alpha>0,\mu\ge0$, let $(u_\infty(\cdot;\alpha,\mu),m_\infty(\cdot;\alpha,\mu))$ be the solution of \eqref{eqn: SS dispersal} that satisfies \eqref{mass conservation}. 
\begin{enumerate}
\item[(i)] There exist constants $\underbar{u}(\alpha),\bar u(\alpha) >0$ such that $ \underbar u(\alpha)\le u_\infty(x;\alpha,\mu)\le \bar u(\alpha)$ for all $x\in \Omega$ and $\mu\ge 0$.
\item[(ii)]  $0<m_\infty(x;\alpha,\mu)\le \alpha\bar s(x)$ for all $\mu\ge 0$. Furthermore, $m_\infty(\cdot;\alpha,\mu)\to 0$ uniformly as $\mu\to\infty$.
    % \item [(iv)] There is a positive-measure subset $N\subset(-\ell,\ell)$ of $m_\infty(x;\mu_1)>m_\infty(x;\mu_2)$ if $\mu_1<\mu_2$ for $x\in N$. 
    % {\color{red}If $\gamma(m)m$ is increasing in $m\ge 0$, $N=(-\ell,\ell)$.}
    % \item [(iv)] There is no subinterval $I\subset (-\ell,\ell)$ such that $\partial _\mu u_\infty(x;\alpha,\mu)= 0$ for all $x\in I$.

    \item [(iii)] There exist constants $C>0$ and $\mu_0$ such that for all $\mu\ge \mu_0$, $u_\infty$ has its maximum at $x=0$, and
\[
    \frac{\partial u_\infty}{\partial \mu}(0;\alpha,\mu)<0 \text{ and } \bigl|\frac{\partial u_\infty}{\partial \mu}(x;\alpha,\mu)\bigr| \;\le\; \frac{C}{\mu^2}
    \qquad \text{for all } x\in[-\ell,\ell].
\]
In particular, $u_\infty(\cdot;\alpha,\mu)\to \frac{M}{2\ell}$ uniformly as $\mu\to\infty$. 
    
    \item[(iv)] If $z\mapsto -z\gamma'(z)$ is increasing on $(0,z_0]$ and decreasing on $(z_0,\infty)$, then there exists $\alpha_1<\alpha_2$ such that $\frac{\partial u_\infty}{\partial \mu}(0;\alpha,0)<0$ for $\alpha <\alpha_1$ and $\frac{\partial u_\infty}{\partial \mu}(0;\alpha,0)>0$ for $\alpha >\alpha_2$.
\end{enumerate}
\end{theorem}

\begin{proof}
\textbf{(i)} 
Recall from \eqref{mass conservation} that the constant $M=\int_\Omega u_\infty(x;\alpha,\mu)dx$ is independent of $\alpha$ and $\mu$. Let $x\in(-\ell,\ell)$ and $\alpha,\mu>0$ be fixed. 
By using the boundary condition again, there exists $K=K(\alpha,\mu)>0$ independent of $x$ such that \begin{equation}\label{eqn:SSrel}
     \gamma(m_\infty)u_\infty=K.
 \end{equation} 
   By using \eqref{eqn:SSrel}, the constant $K$ can be expressed by
\[K = \frac{M}{\int_\Omega\frac{1}{\gamma(m_\infty(x;\alpha,\mu))}dx}.\] Since $\gamma$ is bounded below by $\gamma(\alpha\max  \bar s)$ and bounded above by $\gamma(0)$, the map $\mu\mapsto K(\alpha,\mu)$ is bounded, that is, \[ \frac{\gamma(\max{\alpha \bar s}) M}{|\Omega|}\le K(\alpha,\mu)\le \frac{\gamma(0) M}{|\Omega|}.\] The bounds for $K$, $\gamma$ and \eqref{eqn:SSrel} ensures that
\[\frac{\gamma(\max{\alpha \bar s}) M}{\gamma(0)|\Omega|}\le u_\infty(x;\alpha,\mu)\le \frac{\gamma(0) M}{\gamma(\alpha \max{\bar s})|\Omega|}. \]

 
\textbf{(ii)} The equality $m_\infty = \alpha\bar s\dfrac{u_\infty}{u_\infty +\mu}$ directly yields the bound. Since $u_\infty$ is uniformly bounded in $\mu$, $m_\infty(\cdot;\alpha,\mu) \to 0$ uniformly in $\overline{\Omega}$ as $\mu\to \infty$.

% We now claim that $m(x;\mu)$ is decreasing in $\mu$. For notational convenience, let us omit $x$ for associated functions. Suppose that $m_\infty(\mu_1)\le m_\infty(\mu_2)$ for $\mu_1<\mu_2$. Then $\mu_1m_\infty(\mu_1)<\mu_2 m_\infty(\mu_2)$, so that we have $$(\alpha \bar s-m_\infty(\mu_1))u_\infty(\mu_1)<(\alpha \bar s-m_\infty(\mu_2))u_\infty(\mu_2).$$
% Rearranging gives
% \begin{align*}
%     \alpha \bar s(u_\infty(\mu_1)-u_\infty(\mu_2))&<m_\infty(\mu_1)u_\infty(\mu_1)-m_\infty(\mu_2)u_\infty(\mu_2)\\
%     &\le m_\infty(\mu_2)(u_\infty(\mu_1)-u_\infty(\mu_2)),
% \end{align*} 
% hence by the second statement
% \begin{equation*}
%     u_\infty(\mu_1)-u_\infty(\mu_2)<0.
% \end{equation*}
% Then we have obtained that
% \begin{align*}
%     u_\infty(x;\mu_1)-u_\infty(x;\mu_2)<0 \text{ for all } x\in \Omega.
% \end{align*} If it satisfies for all $x\in[-\ell,\ell]$, it contradicts to the mass conservation of $u$ by (i), which concludes the third statement.

% \textbf{(iii)} We disclaim that $\frac{\partial u_\infty}{\partial \mu}(x;\mu) = 0$ for all $x$ by leading a contradiction. Assume that $\frac{\partial u_\infty}{\partial \mu}(x;\mu) = 0$ for all $x\in\Omega$. 
% From the relation
% \[ u_\infty = \frac{\mu m_\infty}{\alpha \bar s-m_\infty},\] we have
% \begin{align*}
%     &\frac{\partial u_\infty}{\partial \mu} = \frac{m_\infty(\alpha \bar s-m_\infty) + \mu \alpha \bar s\partial_\mu m_\infty}{(\alpha \bar s-m_\infty)^2}=0\\
%     &\iff \mu \partial_\mu m_\infty = \frac{m_\infty^2- \alpha \bar sm_\infty}{\alpha \bar s}.
% \end{align*} By using the separation of variables, we obtain the relation
% \begin{equation*}
%     m_\infty(x;\mu) = \frac{\alpha \bar s(x)}{1+\mu e^{C(x)}}
% \end{equation*} for some function $C(x)$ independent of $\mu$. By using the equation $(\alpha \bar s-m)u-\mu m=0$, we have 
% \begin{equation*}
%     u_\infty(x) = e^{-C(x)}.
% \end{equation*}
% If we plug it into \eqref{eqn:SSrel} and differentiate both sides with respect to $x$, we obtain
% \begin{equation*}
%     C'(x) = \frac{\gamma(m_\infty)_x}{\gamma(m_\infty)}.
% \end{equation*} By integrating it over $(0,x)$, we have
% \begin{equation*}
%     C(x) = C(0) + \ln (\gamma(m_\infty)),
% \end{equation*} which contradicts to the independency of $C(x)$ on $\mu$.

\textbf{(iii)} Recall the constant $K=K(\alpha,\mu)$ from \eqref{eqn:SSrel}. Let $F(u,K,\alpha,\mu) = \gamma(m)u - K$, so we have $F(u_\infty,K,\alpha,\mu)=0$ in $\overline{\Omega}$. Differentiate both sides with respect to $\mu$:
\begin{align}\label{eqn: total deriv}
    \frac{\partial F}{\partial u}  \frac{\partial u_\infty}{\partial \mu} +\frac{\partial F}{\partial K} \frac{\partial K}{\partial \mu}  +\frac{\partial F}{\partial \mu}  =0,
\end{align} where
\begin{align*}
    \frac{\partial F}{\partial u}  = \gamma(m_\infty) +\alpha \bar s\gamma'(m_\infty) \frac{\mu u_\infty}{(u_\infty+\mu)^2}, \,\,\frac{\partial F}{\partial K}  = -1, \text{ and } \frac{\partial F}{\partial \mu}  = -\alpha\bar s\gamma'(m_\infty) \frac{u_\infty^2}{(u_\infty+\mu)^2}.
\end{align*}
Note that $m_\infty$ and $u_\infty$ are uniformly bounded on $\overline{\Omega}$, and $m_\infty \to 0$ as $\mu\to\infty$ by (i) and (ii). Since $\gamma\in C^2[0,\infty)$, $\gamma(m_\infty)\to\gamma(0)$ and $\frac{\partial F}{\partial u} (x)\to \gamma(0)$ uniformly as $\mu\to\infty$.  Choose a sufficiently large $\mu_0>0$ such that for some $c_0>0$,
\[\frac{\partial F}{\partial u}  = \gamma(m_\infty) +\alpha \bar s\gamma'(m_\infty) \frac{\mu u_\infty}{(u_\infty+\mu)^2}>c_0,\quad \frac{\partial F}{\partial u} =\gamma(0)+O(\mu^{-1}), \text{ and, }\]\[\frac{\partial F}{\partial \mu}  = -\alpha\bar s\gamma'(m_\infty)\frac{u_\infty^2}{(u_\infty+\mu)^2}\le \frac{C}{\mu^2} \] for all  $ \mu>\mu_0$, where $C$ is a positive constant independent of $\mu$. We  differentiate both sides of $\gamma(m_\infty)u_\infty=K$ with respect to $x$ to have
\[ \frac{\partial F}{\partial u}u' = -\alpha\frac{u_\infty}{u_\infty+\mu}\gamma'(m_\infty)\bar s'. \] Since $\frac{\partial F}{\partial u}>0$ and $ \gamma'(m_\infty)<0$, $u'$ and $\bar s'$ have the same sign. Therefore, $u$ has its maximum at $x=0$.

The mass constraint from (i) implies
\[
0 = \frac{d}{d\mu}\int_{-\ell}^{\ell} u_\infty\,dx
= \int_{-\ell}^{\ell} \frac{\partial u_\infty}{\partial \mu}\,dx
= \int_{-\ell}^{\ell} \frac{\frac{\partial K}{\partial \mu}  - \frac{\partial F}{\partial \mu} }{\frac{\partial F}{\partial u} }\,dx.
\]
Solving for $\frac{\partial K}{\partial \mu} $ gives
\[
\frac{\partial K}{\partial \mu} 
=
\frac{\displaystyle\int_{-\ell}^{\ell} \frac{\frac{\partial F}{\partial \mu} }{\frac{\partial F}{\partial u} }\,dx}
{\displaystyle\int_{-\ell}^{\ell} \frac{1}{\frac{\partial F}{\partial u} }\,dx}.
\]
The mass constraint $M=\int_{-\ell}^{\ell}u_\infty dx$ and \eqref{eqn:SSrel} gives
\[u_\infty = \frac{K(\mu)}{\gamma(m_\infty)} = \displaystyle\frac{M}{\gamma(m_\infty)\int_{-\ell}^{\ell}\frac{1}{\gamma(m_\infty)}dx}\to \frac{M}{\gamma(0)\frac{2\ell}{\gamma(0)}}=\frac{M}{2\ell}\] uniformly in $x$.

We now claim that $\frac{\partial u_\infty}{\partial \mu}<0$ for sufficiently large $\mu$. For large $\mu$, 
\[\frac{\partial F}{\partial \mu}  = -\alpha\bar s\gamma'(m_\infty) \frac{u_\infty^2}{(u_\infty+\mu)^2}=-\frac{\alpha\bar s(x)\gamma'(0)M}{2\ell\mu^2}+O(\mu^{-2})\] uniformly in $\overline{\Omega}$. Therefore, for large $\mu$
\[\frac{\partial K}{\partial \mu}  = \frac{\displaystyle\int_{-\ell}^{\ell}\frac{\frac{\partial F}{\partial \mu} }{\frac{\partial F}{\partial u} }dx}{\displaystyle\int_{-\ell}^{\ell}\frac{1}{\frac{\partial F}{\partial u} }dx}=\frac{\displaystyle\int_{-\ell}^{\ell}\frac{\frac{\partial F}{\partial \mu} }{\gamma(0)}dx}{\displaystyle\int_{-\ell}^{\ell}\frac{1}{\gamma(0)}dx} +O(\mu^{-1})= \overline{\frac{\partial F}{\partial \mu} } + O(\mu^{-1}),\] where $\overline{\frac{\partial F}{\partial \mu} } = \frac{1}{2\ell}\int_{-\ell}^\ell \frac{\partial F}{\partial \mu} (x)dx$. Solve \eqref{eqn: total deriv} for $\frac{\partial u_\infty}{\partial \mu}$, and we obtain
\[\frac{\partial u_\infty}{\partial \mu} = \frac{\frac{\partial K}{\partial \mu}  - \frac{\partial F}{\partial \mu} (x)}{\frac{\partial F}{\partial u} } = \frac{\overline{\frac{\partial F}{\partial \mu} }-\frac{\partial F}{\partial \mu} (x)}{\gamma(0)}+O(\mu^{-2}).  \]
Since $u_\infty$ and $\bar s$ has its maximum at $x=0$, $\overline{\frac{\partial F}{\partial \mu} }<\frac{\partial F}{\partial \mu} (x)$ for large $\mu$, hence $\frac{\partial u_\infty}{\partial \mu}<0$ for sufficiently large $\mu$. 

Using $\frac{\partial F}{\partial u} >0$ and $|\frac{\partial F}{\partial \mu} |\le C/\mu^2$, we obtain
\[
|\frac{\partial K}{\partial \mu} |
\le
\frac{\displaystyle\int_{-\ell}^{\ell} \frac{|\frac{\partial F}{\partial \mu} |}{\frac{\partial F}{\partial u} }\,dx}
{\displaystyle\int_{-\ell}^{\ell} \frac{1}{\frac{\partial F}{\partial u} }\,dx}
\le
\frac{\displaystyle\frac{C}{\mu^2}\int_{-\ell}^{\ell} \frac{1}{\frac{\partial F}{\partial u} }\,dx}
{\displaystyle\int_{-\ell}^{\ell} \frac{1}{\frac{\partial F}{\partial u} }\,dx}
= \frac{C}{\mu^2}.
\]

Finally, substituting it back into the expression for $\partial_\mu u$,
\[
|\partial_\mu u|
\le
\frac{|\frac{\partial K}{\partial \mu} | + |\frac{\partial F}{\partial \mu} (x,\mu)|}{\frac{\partial F}{\partial u} (x,\mu)}
\le
\frac{1}{c_0}\left(\frac{C}{\mu^2} + \frac{C}{\mu^2}\right)
= \frac{2C}{\mu^2},
\]
for some constant $C>0$ independent of $\mu$ (for all $\mu\ge \mu_0$). This proves the claimed $O(\mu^{-2})$ decay and uniform convergence $\partial_\mu u(\cdot,\mu)\to 0$ as $\mu\to\infty$.

\textbf{(iv)}  At $\mu=0$, we have 
\begin{equation*}
    \frac{\partial F}{\partial u}  = \gamma(\alpha \bar s)>0, \text{ and } \frac{\partial F}{\partial \mu}  = -\alpha\bar s\gamma'(\alpha\bar s)>0.
\end{equation*} Therefore, $u$ has its maximum at $x=0$ by the relation $\gamma(\alpha\bar s)u' = -\alpha\gamma'(\alpha \bar s)\bar s'$.  
If we combine the constraint and \eqref{eqn: total deriv}, we obtain 
\begin{align*}
    \frac{\partial K}{\partial \mu} \Big|_{{\mu=0}} = \frac{\int_{-\ell}^{\ell}\frac{\partial F}{\partial \mu}  \frac{\partial F}{\partial u} ^{-1}dx}{\int_{-\ell}^{\ell}\frac{\partial F}{\partial u} ^{-1}dx} =\frac{\int_{-\ell}^{\ell}-\alpha\bar s\gamma'(\alpha\bar s) \gamma(\alpha \bar s)^{-1}dx}{\int_{-\ell}^{\ell}\gamma(\alpha \bar s)^{-1}dx}.
\end{align*} Notice that $\frac{\partial K}{\partial \mu} $ is a weighted average of $-\alpha\bar s \gamma'(\alpha \bar s)$ over $(-\ell,\ell)$. We use \eqref{eqn: total deriv} again to have
\begin{equation*}
    \partial_\mu u\Big|_{\mu=0} = \frac{\frac{\partial K}{\partial \mu} \big|_{\mu=0} -(-\alpha\bar s\gamma'(\alpha\bar s))}{\gamma(\alpha\bar s)},
\end{equation*} which implies that the sign of $\partial_\mu u\big|_{\mu=0}$ is determined by the sign of $\frac{\partial K}{\partial \mu} \big|_{\mu=0} -(-\alpha\bar s\gamma'(\alpha\bar s))$. Suppose that the map $z\mapsto -z\gamma'(z)$ is increasing on $(0,z_0)$ and decreasing on $(z_0,\infty)$. Choose $\alpha_1 = z_0/\bar s(0)$ so that  $(\alpha\bar s(\ell),\alpha\bar  s(0))\subset (0,z_0)$ for $\alpha<\alpha^*$. Since $\bar s(0)$ is the maximum of $\bar s$, for $\alpha<\alpha_1$
\begin{equation*}
    -\alpha\bar s(x)\gamma'(\alpha\bar s(x)) - (-\alpha\bar s(0)\gamma'(\alpha\bar s(0))),
\end{equation*} hence $\frac{\partial K}{\partial \mu} \big|_{\mu=0} -(-\alpha\bar s\gamma'(\alpha\bar s))<0$, so $\partial_\mu u\big|_{\mu=0}<0$ at $x=0$. Let $\alpha_2 = z_0/\bar s(\ell)$. Then $(\alpha \bar s (\ell),\alpha \bar s (0))\subset (z_0,\infty)$. Similarly, for $\alpha>\alpha_2$, we have $\partial_\mu u\big|_{\mu=0}>0$ at $x=0$.
% We now prove (vi). First, we claim that $\mu\mapsto u_\infty(x;\mu)/\mu$ is decreasing in $\mu>0$ for each $x\in\Omega$. We have
% \[\frac{\mu}{u_\infty+\mu} = \frac{\mu}{\frac{\mu m_\infty}{\alpha\bar s-m_\infty}+\mu}=\frac{\alpha\bar s -m_\infty}{\alpha \bar s}. \] By differentiating both sides with respect to $\mu$, we get
% \[\partial_\mu \left(\frac{\mu}{u_\infty+\mu}\right)=-\frac{\partial_\mu m_\infty}{\alpha \bar s}>0,\] where we have used (iv). Expanding the derivative and rearranging the equation yields
% \[ \frac{u_\infty-\mu\frac{\partial u_\infty}{\partial \mu}}{(u+\mu)^2}>0 \implies u_\infty-\mu\frac{\partial u_\infty}{\partial \mu}>0 \text{ for all } x\in\Omega \text{ and } \mu>0. \]
% Then \[ \partial_\mu\left(\frac{u_\infty}{\mu}\right) = \frac{\mu \frac{\partial u_\infty}{\partial \mu}-u_\infty}{\mu^2}<0.\] Therefore, for each $t\in(0,1]$, \[ u_\infty(t\mu) = t\mu\frac{u_\infty(t\mu)}{t\mu}\ge t\mu\frac{u_\infty(\mu)}{\mu} = tu_\infty(\mu) \] for all $\mu>0$. For any $(x_0,y_0)$ such that $y_0\le u(x_0)$ and $t\in[0,1]$, we have  $ty_0\le tu(x_0)\le u(tx_0)$. Therefore, the positive region under the graph $(\mu,u_\infty(\mu))$ is star-shaped centered at the origin.

\end{proof} 

\end{document}
